\documentclass[11pt,aspectratio=169]{beamer}

\usetheme{Madrid}
\usecolortheme{beaver}
\setbeamertemplate{navigation symbols}{}

\usepackage[utf8]{inputenc}
\usepackage[T1]{fontenc}
\usepackage{amsmath,amssymb}
\usepackage{booktabs}
\usepackage{listings}
\usepackage{xcolor}
\usepackage{tikz}
\usepackage{tcolorbox}
\tcbuselibrary{skins,breakable}

\definecolor{codebg}{RGB}{245,245,245}
\definecolor{codekw}{RGB}{0,0,180}
\lstdefinestyle{latexcode}{
  backgroundcolor=\color{codebg},
  basicstyle=\ttfamily\scriptsize,
  keywordstyle=\color{codekw}\bfseries,
  commentstyle=\color{gray},
  frame=single,
  rulecolor=\color{gray!40},
  breaklines=true,
  showstringspaces=false,
  xleftmargin=4pt,
  xrightmargin=4pt,
}
\lstset{style=latexcode}

\newtcolorbox{outputbox}{
  colback=white, colframe=black!65, boxrule=0.5pt, arc=1pt,
  left=8pt, right=8pt, top=5pt, bottom=4pt,
  title=\textsc{Compiled Output}, fonttitle=\bfseries\scriptsize,
  coltitle=white, colbacktitle=black!65
}
\newcommand{\whybox}[1]{%
  \vspace{2pt}
  {\scriptsize\raggedright\textcolor{red!55!black}{\textbf{Why \& What:}}~#1\par}
}

\newtcolorbox{exercisebox}{
  colback=orange!8, colframe=orange!70!black, boxrule=0.6pt,
  arc=2pt, left=6pt, right=6pt, top=4pt, bottom=4pt,
  title={Hands-on Exercise}, fonttitle=\bfseries, coltitle=white,
  colbacktitle=orange!70!black
}
\newtcolorbox{homeworkbox}{
  colback=blue!6, colframe=blue!50!black, boxrule=0.6pt,
  arc=2pt, left=6pt, right=6pt, top=4pt, bottom=4pt,
  title={Homework}, fonttitle=\bfseries, coltitle=white,
  colbacktitle=blue!50!black
}

\title[\texttt{\textbackslash begin} to \texttt{\textbackslash end}]{From \texttt{\textbackslash begin} to \texttt{\textbackslash end}}
\subtitle{The Complete LaTeX Guide \\[2pt] \normalsize Day 5 of 7: Document Structure \& Long-Document Tools}
\author[A.\ Rahman]{Atikur Rahman}
\institute[White Globe Initiative]{White Globe Initiative}

\setbeamertemplate{footline}{
  \leavevmode
  \hbox{%
  \begin{beamercolorbox}[wd=.38\paperwidth,ht=2.5ex,dp=1.125ex,leftskip=1em]{author in head/foot}%
    \usebeamerfont{author in head/foot}White Globe Initiative
  \end{beamercolorbox}%
  \begin{beamercolorbox}[wd=.42\paperwidth,ht=2.5ex,dp=1.125ex,center]{title in head/foot}%
    \usebeamerfont{title in head/foot}From \texttt{\textbackslash begin} to \texttt{\textbackslash end} --- Day 5
  \end{beamercolorbox}%
  \begin{beamercolorbox}[wd=.20\paperwidth,ht=2.5ex,dp=1.125ex,right,rightskip=1em]{date in head/foot}%
    \usebeamerfont{date in head/foot}\insertframenumber{} / \inserttotalframenumber
  \end{beamercolorbox}}%
  \vskip0pt
}

\begin{document}

\begin{frame}[plain]
  \begin{center}
    \vspace{4pt}
    {\Huge\bfseries\usebeamercolor[fg]{title}From \texttt{\textbackslash begin} to \texttt{\textbackslash end}}\\[6pt]
    {\Large\itshape The Complete LaTeX Guide}\\[18pt]
    {\large\bfseries Day 5 of 7 --- Document Structure \& Long-Document Tools}\\[22pt]
    \rule{0.7\textwidth}{0.4pt}\\[14pt]
    \begin{tabular}{rl}
      \textbf{Instructor:}       & Atikur Rahman \\[3pt]
      \textbf{Organized by:}     & White Globe Initiative \\[3pt]
      \textbf{Instructor Email:} & \texttt{atik@whiteglobeinitiative.org} \\[3pt]
      \textbf{General Contact:}  & \texttt{contact@whiteglobeinitiative.org} \\[3pt]
      \textbf{Format:}           & 7-Day Intensive Course \\[3pt]
      \textbf{Session:}          & Day 5 of 7 \\
    \end{tabular}\\[18pt]
    \rule{0.7\textwidth}{0.4pt}
  \end{center}
\end{frame}

\begin{frame}{Today's Agenda}
  \tableofcontents
\end{frame}

\section{Recap \& Today's Goal}

\begin{frame}{Quick Recap: Day 4}
  \small
  Yesterday: BibTeX \texttt{.bib} files, \texttt{natbib}, author-year
  citation styles, \texttt{citet} vs.\ \texttt{citep}, and
  cross-referencing sections with \texttt{\textbackslash
  label}/\texttt{\textbackslash ref}.
  \vspace{6pt}

  \textbf{Today's goal:} move from a single short file to the
  structure a full paper or thesis-length document needs, with the
  navigation aids a reader -- or a referee -- expects.
\end{frame}

\section{Table of Contents}

\begin{frame}[fragile]{Table of Contents}
  \begin{lstlisting}
\begin{document}
\maketitle
\tableofcontents

\section{Introduction}
\section{Data and Methodology}
\section{Results}
  \end{lstlisting}
  \begin{outputbox}
    \small
    \textbf{Contents}\\[4pt]
    1\quad Introduction \dotfill 1\\
    2\quad Data and Methodology \dotfill 2\\
    3\quad Results \dotfill 4
  \end{outputbox}
  \whybox{\texttt{\textbackslash tableofcontents} auto-generates a full
  TOC from every \texttt{\textbackslash section} in your document --
  you type nothing else. It \textbf{updates automatically} on every
  recompile, including page numbers, so it never goes stale as you
  edit. It needs \textbf{two} compiles to settle those page numbers
  correctly the first time (Overleaf does the second pass for you).}
\end{frame}

\section{Footnotes}

\begin{frame}[fragile]{How to Include a Footnote}
  \begin{lstlisting}
Inflation targeting has been the dominant framework
since the 1990s.\footnote{See \citet{nandi2024crude}
for a related discussion of exchange-rate volatility.}
  \end{lstlisting}
  \begin{outputbox}
    \small Inflation targeting has been the dominant framework since
    the 1990s.\textsuperscript{1}\\[14pt]
    \rule{3.2cm}{0.3pt}\\
    {\tiny \textsuperscript{1}See Nandi (2024) for a related discussion of exchange-rate volatility.}
  \end{outputbox}
  \whybox{\texttt{\textbackslash footnote\{...\}} does two things at
  once: it drops a small superscript number right where you typed the
  command (see the ``1'' after ``1990s.'' above), and it automatically
  prints the matching numbered note at the \textbf{bottom of that same
  page}, below a short rule -- exactly as shown in the output box.
  Numbering is continuous and automatic across the whole document, so
  you never assign numbers yourself. Economics convention: use
  footnotes for asides or methodological details, and use citations
  (Day 4) for sources -- don't use footnotes as a second bibliography.}
\end{frame}

\section{Hyperlinks}

\begin{frame}[fragile]{hyperref: Clickable Cross-References}
  \begin{lstlisting}
\usepackage{hyperref}   % load LAST in the preamble

See Section~\ref{sec:strategy} or visit
\href{https://whiteglobeinitiative.org}{our website}.
  \end{lstlisting}
  \begin{outputbox}
    \small See Section~\textcolor{blue}{2} or visit
    \textcolor{blue}{\underline{our website}}.
  \end{outputbox}
  \whybox{Loading \texttt{hyperref} turns every \texttt{
  \textbackslash ref}, \texttt{\textbackslash eqref}, \texttt{
  \textbackslash cite}, and TOC entry into a \textbf{clickable link}
  in the PDF -- shown in blue above, though colors are configurable.
  \texttt{\textbackslash href\{url\}\{text\}} adds a real hyperlink.
  \textbf{Load order matters}: put \texttt{hyperref} as the
  \emph{last} \texttt{\textbackslash usepackage} in your preamble to
  avoid conflicts with other packages.}
\end{frame}

\section{Multi-file Projects}

\begin{frame}[fragile]{Splitting a Long Paper Across Files}
  \vspace{-4pt}
  \begin{lstlisting}[basicstyle=\ttfamily\tiny,aboveskip=1pt,belowskip=1pt]
% main.tex
\documentclass{article}
\usepackage{amsmath, natbib, hyperref}
\begin{document}
\maketitle
\tableofcontents
\input{introduction}
\input{literature}
\input{methodology}
\bibliography{references}
\end{document}
  \end{lstlisting}
  \begin{outputbox}
    \scriptsize\ttfamily
    Project files: main.tex, introduction.tex, literature.tex,
    methodology.tex, references.bib\\[3pt]
    \normalfont\footnotesize One combined PDF -- identical to writing
    it all in a single file.
  \end{outputbox}
  \whybox{\texttt{\textbackslash input\{filename\}} pulls in
  \texttt{filename.tex} at that exact point -- write each section in
  its own file. The \textbf{compiled PDF looks no different}; only
  your project became easier to navigate. \texttt{\textbackslash
  include} is a close cousin, better suited to \emph{chapters}.}
\end{frame}

\section{Page Numbering \& Headers}

\begin{frame}[fragile]{Page Numbering, Headers and Footers}
  \vspace{-4pt}
  \begin{lstlisting}[aboveskip=1pt,belowskip=1pt]
\usepackage{fancyhdr}
\pagestyle{fancy}
\rhead{Draft -- Do Not Circulate}
\lhead{Remittances and Consumption Smoothing}
  \end{lstlisting}
  \begin{outputbox}
    \scriptsize
    Remittances and Consumption Smoothing \hfill Draft -- Do Not Circulate\\
    \rule{\linewidth}{0.3pt}\\[6pt]
    \footnotesize (page body here)\\[6pt]
    \rule{\linewidth}{0.3pt}\\
    \hfill 3 \hfill
  \end{outputbox}
  \whybox{Default \texttt{article} numbering is centered at the bottom
  -- fine for most submissions. \texttt{fancyhdr} lets you add running
  headers, like the short title and a draft watermark shown above --
  common for working-paper versions. Check your target journal's
  template first (Day 6): most already configure this for you.}
\end{frame}

\section{Practice}

\begin{frame}{Hands-on Exercise}
  \begin{exercisebox}
    Split your running paper draft into separate files:
    \texttt{introduction.tex}, \texttt{literature.tex},
    \texttt{methodology.tex}, \texttt{results.tex},
    \texttt{conclusion.tex}. Assemble them into a \texttt{main.tex}
    using \texttt{\textbackslash input}, add
    \texttt{\textbackslash tableofcontents}, and add at least one
    \texttt{\textbackslash footnote} somewhere in the text. Confirm the
    document compiles to a single, correctly ordered PDF.
  \end{exercisebox}
\end{frame}

\begin{frame}{Homework Before Day 6}
  \begin{homeworkbox}
    Load \texttt{hyperref} and confirm every cross-reference
    (equations, tables, figures, sections, citations) is clickable and
    resolves correctly -- fix any \texttt{??} you find.
  \end{homeworkbox}
  \vspace{10pt}
  \small\textbf{Resources for today:}
  \begin{itemize}\small
    \item Overleaf -- ``Sections and chapters'' / ``Table of contents''
    \item Overleaf -- ``Footnotes'' / ``Hyperlinks''
    \item Overleaf -- ``Multi-file LaTeX projects''
    \item Overleaf -- ``Page numbering'' / ``Headers and footers''
    \item Wikibooks LaTeX -- Document Structure chapter
  \end{itemize}
\end{frame}

\begin{frame}{Coming Up: Day 6}
  \Large
  \centering
  Journal Templates \& Professional Formatting \\[6pt]
  \normalsize
  Migrating your full paper into a real economics-journal template,
  end to end.
\end{frame}

\begin{frame}[plain]
  \begin{center}
    \vspace{35pt}
    {\Large\bfseries Thank You}\\[10pt]
    {\large From \texttt{\textbackslash begin} to \texttt{\textbackslash end}: The Complete LaTeX Guide}\\[4pt]
    Day 5 of 7 --- Document Structure \& Long-Document Tools\\[20pt]
    \rule{0.5\textwidth}{0.4pt}\\[10pt]
    Atikur Rahman \\
    White Globe Initiative \\[2pt]
    \texttt{atik@whiteglobeinitiative.org} \\
    \texttt{contact@whiteglobeinitiative.org}
  \end{center}
\end{frame}

\end{document}
